%! Inference and Experiments — Unit 3, Lesson 3.5 — Homework
% GRADUAL RELEASE: this is the lesson's SCORED practice and the LAST component in the packet.
% Paper homework is the DEFAULT for this course; DeltaMath is an occasional override that
% replaces this page. Started in the last eight minutes of the period and finished after two
% study halls. The AP Practice component before it is extra credit, scored 1 to 5.
% THIRD CONTEXT: the notes teach on studying with music, Guided Practice runs on a reading app,
% AP Practice on sports drinks, tomatoes, and green tea; these items are on a cafeteria
% "fruit first" experiment and a city walking survey.
% Part A (Fruit First — random selection AND random assignment, a GIVEN simulation dotplot):
%   A1 read the dotplot (3 of 100 at 0.4 or more), A2 significant in context, A3 cause and to whom.
%   Simulated differences (fruit first − usual line), 100 reshuffles:
%   -0.5:1 -0.4:2 -0.3:6 -0.2:11 -0.1:17 0:26 0.1:17 0.2:11 0.3:6 0.4:2 0.5:1
% Part B (Walking and Stress — a random sample, NO random assignment): B1 the AP-style multiple
%   choice with a required justification, B2 THE CRUX restated — "chosen at random, so this
%   proves walking lowers stress". 9 points.
% \answerspace{H}{} reserves exactly H; each \gradedblank inside it is BYTE-IDENTICAL in the key,
% where key mode prints its display. Keep every \boxguard identical in the blank and the key.
\documentclass[12pt]{article}
\usepackage{apstats-article}
\usepackage{apstats-boxes}
\usepackage{grading}
\gradingassignment{apstats-u03-l05-hw}
\gradingcourse{advanced_statistics}
\newcommand{\answerspace}[2]{\par\nopagebreak\noindent\begin{minipage}[t][#1][t]{\linewidth}%
  \color{keyred}\bfseries #2\end{minipage}\par}

\begin{document}

\pageheader{Unit 3, Lesson 3.5}{Homework}
{\small\itshape Write each final answer in its box. Only what is inside a box is graded.}\par
% NAMESTRIP: no \namedateperiod here — the name row lives on the cover only.

\vspace{0.04in}

% Authored BYTE-IDENTICALLY in homework/ and homework_key/.
\boxguard[8]
\begin{remindbox}
\small
\textbf{This is your graded homework.} It is scored 1 to 5 and due the first class after two study
halls. Start it now; ask about anything that stalls you before you leave.
\end{remindbox}

\vspace{0.04in}

\boxguard[30]
\begin{notesbox}{Part A --- Fruit First}
\interlinepenalty=10000 % keep each prompt whole when the box breaks
\small
A nutrition researcher picks \textbf{60 students at random} from the 1{,}200 students at a high
school and randomly assigns 30 to a lunch line with the fruit placed \emph{first} and 30 to the
usual line, with the fruit at the end. The fruit-first group takes \textbf{0.4 more pieces of
fruit} per student, on average. A computer reshuffled the 60 results into two new groups of 30,
100 times, as if the order of the line made no difference.
\begin{center}
\begin{tikzpicture}[x=0.95cm, y=0.1cm]
  \draw[thick] (-5.4,0) -- (5.4,0);
  \foreach \x/\l in {-5/-0.5,-4/-0.4,-3/-0.3,-2/-0.2,-1/-0.1,0/0,1/0.1,2/0.2,3/0.3,4/0.4,5/0.5} {
    \draw (\x,0) -- (\x,-1.4); \node[below, font=\tiny] at (\x,-1.2) {$\l$};}
  \foreach \x/\n in {-5/1,-4/2,-3/6,-2/11,-1/17,0/26,1/17,2/11,3/6,4/2,5/1}
    \foreach \k in {1,...,\n} \fill[navy] (\x,\k) circle (1.3pt);
  \draw[redacc, dashed, thick] (4,3) -- (4,14);
  \node[font=\tiny\bfseries, redacc, above] at (4,14) {observed 0.4};
  \node[font=\tiny] at (0,-7) {Simulated difference in mean pieces of fruit, fruit first $-$ usual};
  \node[font=\scriptsize\bfseries] at (0,30) {100 random reshuffles};
\end{tikzpicture}
\end{center}
\begin{enumerate}[label=\textbf{A\arabic*.}, leftmargin=2.2em, itemsep=2pt]
  \item How many of the 100 reshuffles gave a difference of 0.4 pieces or more?
        \hfill\rule[-0.17in]{0pt}{0.43in}\smash{\gradedblank[id=hw1, type=short, points=1, answer=3]{1.6cm}{0.3in}}
  \item Is the observed difference of 0.4 pieces of fruit statistically significant? Explain in
        context.
        \answerspace{2.1in}{\normalcolor\mdseries\hspace*{10.5pt}\gradedblank[id=hw2, type=rubric, points=2, display={Yes. Only 3 of 100 reshuffles gave 0.4 or more, so a difference this large would rarely happen by chance alone if the order of the line did not matter.}]{\linewidth-10.5pt}{0.95cm}}
  \item Can the researcher conclude that putting fruit first \emph{caused} students to take more
        fruit? To whom can she generalize the result? Explain both.
        \answerspace{2.1in}{\normalcolor\mdseries\hspace*{10.5pt}\gradedblank[id=hw3, type=rubric, points=2, display={Yes: the lines were randomly assigned and the difference is significant. Generalize to all 1,200 students at this school, because the 60 were picked at random from them.}]{\linewidth-10.5pt}{0.95cm}}
\end{enumerate}
\end{notesbox}

\vspace{0.04in}

\boxguard[30]
\begin{notesbox}{Part B --- Walking and Stress}
\interlinepenalty=10000 % keep each prompt whole when the box breaks
\small
A city health department picks \textbf{400 adult residents at random} from a list of all adult
residents. Residents who say they walk at least 30 minutes a day report lower stress, on
average, than residents who do not, and the difference is statistically significant.
\begin{enumerate}[label=\textbf{B\arabic*.}, leftmargin=2.2em, itemsep=2pt]
  \item \textbf{AP-style.} Which conclusion is most appropriate? Write its letter in the Letter
        box, then justify it in a sentence.
        \begin{enumerate}[label=\textbf{(\Alph*)}, leftmargin=2.2em, itemsep=0pt]
          \item Walking causes lower stress for all adults in the city.
          \item Walking causes lower stress, but only for these 400 residents.
          \item Among all adults in the city, walkers tend to report lower stress, but we cannot
                conclude that walking causes it.
          \item No conclusion is possible, because this was not an experiment.
          \item Walking causes lower stress for all adults in the country.
                \hfill\gradedright{\textit{Letter:}~\gradedblank[id=hw4, type=choice, points=1, answer=C]{0.45in}{0.24in}}\rule[-0.2in]{0pt}{0.48in}
        \end{enumerate}
        \answerspace{2.1in}{\normalcolor\mdseries\hspace*{10.5pt}\gradedblank[id=hw4j, type=rubric, points=1, display={(C): the random sample lets us generalize the association to the city's adults, but no one was assigned to walk, so we cannot conclude cause.}]{\linewidth-10.5pt}{0.7cm}}
  \item A council member says, ``The 400 residents were chosen at random, so this proves that
        walking lowers stress.'' Respond, and name a variable that could explain the result
        instead.
        \answerspace{2.1in}{\normalcolor\mdseries\hspace*{10.5pt}\gradedblank[id=hw5, type=rubric, points=2, display={Disagree. Random selection lets us generalize, but only random assignment shows cause. Residents chose whether to walk; walkers may have more free time or less stressful jobs.}]{\linewidth-10.5pt}{0.95cm}}
\end{enumerate}
\end{notesbox}

\vspace{0.04in}

% The next-lesson preview closes the packet, so it lives here rather than in AP Practice.
% A one-line, unbreakable box: guard it by its own height so it moves whole.
\boxguard[4]
\begin{spiralbox}
\small
\textbf{Next: Lesson 4.1, Random and Non-Random Patterns?} Today a simulation told us what
chance alone looks like. Unit 4 opens on that question: is a pattern real, or just chance?
\end{spiralbox}

\end{document}
